\section{新生产关系：供给侧能力、链接网络与组织重写}

上一章讲模型和产品，本章讲生产关系。吴明辉认为，AI 会对软件行业、企业级服务和产业互联网产生重构，生成新的链接关系和生产关系。特别危险的是只提供链接网络、但不掌握供给侧能力、且网络又不是根节点的公司。因为 AI 会降低连接和信息处理成本，削弱中间层的价值。

\begin{figure}[H]
\centering
\includegraphics[width=0.94\textwidth]{supply-demand-relation.png}
\caption{供给侧能力 vs 链接网络：只做连接但不掌握供给侧能力的网络会变危险。自制概念图，依据 02:31:36--02:34:28 对谈内容整理。}
\end{figure}

\begin{knowledgebox}{读图：AI 会重新定价中间层}
如果一家公司只做撮合、分发或关系链，却没有产品、数据、工具、模型或交付能力，当 Agent 能直接连接需求和供给时，它的价值会被压缩。供给侧能力是能真正生产结果的能力。
\end{knowledgebox}

\subsection{新的生产关系：人、模型、工具和数据重新分工}

上一节讲中间层如何被重新定价，本节进一步看工作本身如何被拆开。AI 带来的新生产关系不是“AI 替代所有人”，而是把人、模型、工具、数据和组织重新分工。模型负责认知，工具负责执行接口，数据提供上下文，Agent 负责自动协作，人负责目标、判断和责任。企业服务的未来不是只有 API，而是围绕任务结果重新组织生产。

\begin{figure}[H]
\centering
\includegraphics[width=0.92\textwidth]{new-production-relation.png}
\caption{AI 带来的新生产关系：模型、工具、Agent、私有数据和人类合伙人重新分工。自制概念图，依据 00:00:42--00:00:55 与 02:31:36--02:34:28 对谈内容整理。}
\end{figure}

\begin{importantbox}{生产关系变化的判断标准}
如果 AI 只是提高单点效率，它是工具升级；如果 AI 改变谁拥有数据、谁执行任务、谁承担责任、谁获得收益，它就是生产关系变化。
\end{importantbox}

\begin{knowledgebox}{术语消化：链接、供给侧与生产关系}
\begin{tabularx}{\textwidth}{p{0.22\textwidth}p{0.32\textwidth}X}
\toprule
术语 & 在访谈中的意思 & AI 时代的变化 \\
\midrule
链接网络 & 连接客户、商家、信息或服务的网络 & Agent 降低连接成本，中间层价值被压缩。 \\
供给侧能力 & 真正生产商品、服务、模型或交付结果的能力 & 越靠近结果，越不容易被绕过。 \\
生产关系 & 谁组织生产、谁执行、谁分配收益和责任 & 模型与 Agent 加入后，角色重新分配。 \\
合伙人 & 能承担结果、定义问题、协同 Agent 的人 & 重复执行岗位减少，高责任角色增多。 \\
\bottomrule
\end{tabularx}
\end{knowledgebox}

\subsection{老板与合伙人：组织形态的想象}

访谈中还有一个极端判断：将来企业可能只有两类人，老板和合伙人。这里的合伙人不是传统员工，而是能与 AI、Agent 和工具共同承担结果的人。这个判断未必会字面成真，但它表达了一个趋势：当执行被 AI 自动化，组织会更强调目标定义、资源整合、责任承担和高密度协作。

\begin{figure}[H]
\centering
\includegraphics[width=0.94\textwidth]{boss-partner-org.png}
\caption{老板与合伙人组织：企业可能从雇员组织转向少数老板和大量 AI 合伙人。自制概念图，依据 02:53:20--03:16:40 对谈内容整理。}
\end{figure}

\begin{warningbox}{不要把“只有老板和合伙人”理解成组织立即消失}
现实企业仍需要岗位、流程、合规和管理。但长期看，重复执行型岗位会被 AI 压缩，能定义问题、整合资源、承担结果的人会更像合伙人。
\end{warningbox}

\subsection{本章小结}

AI 对企业服务的冲击不止是效率提升。它会重写链接网络、供给侧能力、组织结构和责任分配。真正危险的是只有中间层价值、没有生产能力的公司。

